\subsection{The characteristic-prime propagation clause in rank four}
\label{subsec:rh-four-p1}

The original \texttt{p1} clause can be obtained from the
ordinary \texttt{P1core} theorem already proved in all ranks.
The essential point is the regularity condition on the
cyclotomic kernel of a rank-four characteristic-prime curve.
We give a direct proof of this condition, using the connected
components of a linear algebraic group and a degree-four
Newton polygon calculation.

\begin{lemma}[A regular spectrum in each connected coset]
\label{lem:rh-four-regular-coset}
Let $H$ be a Zariski-closed subgroup of $\mathrm{GL}_4$ over
an algebraically closed field $E$ of odd characteristic.
Suppose $H$ acts irreducibly, its identity component $H^\circ$
does not act as a sum of four characters, and $H^\circ$ contains
a torus whose four weights on the given representation are
distinct. Then every coset of $H^\circ$ in $H$ contains a
matrix with four distinct eigenvalues.
\end{lemma}
\begin{proof}
Restriction to the normal subgroup $H^\circ$ is semisimple:
choose a simple invariant subspace and take the sum of its
$H$-translates. This sum is semisimple, is nonzero and is all
of $E^4$. No isomorphic simple summand is repeated since that
would repeat a weight of the given torus. The action of $H$
on these distinct summands is transitive, by irreducibility.
Thus either $H^\circ$ acts irreducibly, or there are precisely
two summands of dimension two.

Choose a cocharacter of the torus with four distinct integer
weights $w_1>w_2>w_3>w_4$, after permutation. Such a
cocharacter exists by avoiding finitely many proper rational
hyperplanes in the cocharacter lattice. Work in its eigenbasis.
Consider first a coset which preserves all the simple summands.
On this coset each diagonal-entry function and each diagonal
entry of the inverse is nonzero somewhere: on each summand
the linear span of the matrices of $H^\circ$ is its full
matrix algebra by Burnside's theorem, and this remains true
after multiplication by the fixed coset representative.
The inverse coset has the same property. The coset is an
irreducible variety; choose a matrix $M$ where all these
functions are nonzero. In particular the first diagonal entry
and the last diagonal cofactor of $M$ are nonzero.

Multiply $M$ by $\operatorname{diag}(s^{w_1},\ldots,s^{w_4})$.
The first, third and fourth elementary coefficients of its
characteristic polynomial have degrees at infinity
$w_1$, $w_1+w_2+w_3$, and $\sum_i w_i$, respectively.
For the second coefficient the largest possible degree is
$w_1+w_2$, occurring only for the principal minor indexed
by $1,2$. If this term is nonzero, the four slopes of the
Newton polygon are distinct. Otherwise the degree is at most
$w_1+w_3$, strictly below
$w_1+(w_2+w_3)/2$. The upper Newton polygon then has slopes
$w_1,(w_2+w_3)/2,w_4$ with horizontal lengths $1,2,1$.
For the middle edge its reduced polynomial after the
corresponding scaling is a binomial of degree two with
nonzero constant term. In odd characteristic it is
separable, and the two one-step edges also give simple
roots. In either case the characteristic polynomial over
the Laurent-series field at infinity is separable.
Consequently its discriminant as a Laurent polynomial in
$s$ is not identically zero, and some $s\in E^\times$ gives
a simple spectrum.

For a coset interchanging the two two-dimensional summands,
write its matrices in the form
$\left(\begin{smallmatrix}0&B\\C&0\end{smallmatrix}\right)$.
By the same full-matrix-span argument there is a matrix in
the coset with the $(1,1)$ entries of both off-diagonal blocks
nonzero, after sorting the weights $a_1>a_2$ and $b_1>b_2$
within the respective blocks. The blocks are invertible.
Multiplication by the cocharacter gives squared eigenvalues
with characteristic polynomial
\[
 \det(XI_2-B\operatorname{diag}(s^{b_1},s^{b_2})
                C\operatorname{diag}(s^{a_1},s^{a_2})).
\]
Its trace has degree $a_1+b_1$ with nonzero leading
coefficient, and its determinant has degree
$a_1+a_2+b_1+b_2<2(a_1+b_1)$. Hence the two squared
eigenvalues are distinct and nonzero over $E(s)$.
Their two square roots each are distinct in odd
characteristic. Specializing away from the zeros of the
discriminant gives the assertion also in this case.
\end{proof}

\begin{theorem}[The original rank-four characteristic-prime clause]
\label{thm:rh-four-p1}
Keep the fitted deep tower and the ideal $J$ of
Theorem~\ref{thm:rh-four-image-ideal}. Let
$y:X.A\to k[[z]]$ be any of the original characteristic-prime
points in \texttt{PotentialModularityTower.p1}: the coefficient
residue is the prescribed one, $J$ vanishes, the generic
symplectic representation is absolutely irreducible, and for
every selected local place and every pair of labels the two
characters are distinct on $T.U1(v)$. Then
$X.\mathrm{LocZero}(J,y)$ holds.
\end{theorem}
\begin{proof}
Write $R$ for its generic symplectic representation and
$N=G_L\cap\ker\epsilon$. The coefficient field of the
curve is finite, and the matrices and diagonal characters
are continuous in the $z$-adic topology: this follows from
the original maximal-ideal continuity and the fact that
the coefficient maximal ideal maps into $(z)$.
The same holds for the representation
$\rho=R|_{G_L}\bar\kappa$ obtained through either untwisted
link in characteristic $\ell$. We may use the algebraic
closure $E$ of $k((z))$ as the field of linear algebra.

First $R|_{G_L}$ is irreducible. Otherwise restriction to
the normal subgroup of index two is semisimple by the sum
of translates argument. Its simple constituents can have
at most two isomorphism classes and their multiplicities
are equal. In fact multiplicities cannot exceed one here:
choose an inertia element with four eigenvalues whose
pairwise ratios have infinite order, as constructed below.
Such an element is in $G_L$, and repeated isomorphic
summands would repeat its eigenvalues. The two distinct
constituents must therefore be interchanged by the outside
coset. This gives a self-twist of $R$ by the quadratic
character of $L/F$, using the endomorphism with opposite
signs on the two summands. Reduction of characteristic
polynomials gives a permutation of the original residual
characters by this quadratic character. No two of them
can have the same restriction to $G_L$, by \texttt{Ok.g3}.
Thus such a permutation would fix each character and the
quadratic character would be trivial, a contradiction.
The scalar twist by $\bar\kappa$ does not affect
irreducibility.

To obtain the asserted inertia element, fix one selected
place. On $T.U1(v)$ each of the pairwise ratios of the
graded characters lies in $1+zk[[z]]$, and is nontrivial
by hypothesis. This target is torsion-free: its prime-to-$\ell$
power maps have trivial kernel, and in characteristic
$\ell$ the $\ell$-power map on the domain is injective.
Each ratio has infinite image, so its closed kernel in the
compact local subgroup has infinite index and empty
interior. A finite union of such kernels cannot cover the
group. An element outside their union has every
eigenvalue ratio of infinite order, since the ordinary
flag gives its eigenvalues as the labeled characters.
This element belongs to $N$ by the definition of
$T.U1(v)$. Its powers all have distinct eigenvalues.

Now restrict $\rho$ to $N$. Again the sum of translates
of any simple constituent gives semisimplicity. The
preceding element excludes repeated isomorphic
constituents. If there is more than one constituent, their
number is two or four by transitivity and equality of
dimensions. The quotient $G_L/N$, a subgroup of
$\mathbb Z_\ell^\times$, acts through a transitive
abelian permutation group on these constituents.
A faithful transitive abelian permutation group is
regular, so it has order two or four. It has a nontrivial
quadratic character $\eta$, trivial on a constituent
stabilizer. Multiplication by the value of $\eta$ on a
coset representative, on the corresponding summand,
gives an intertwiner between $\rho$ and its quadratic
twist. The finite permutation action is continuous:
the stabilizers of subspaces are closed, since the
matrices are continuous, and a closed finite-index
subgroup of a compact group is open. Consequently
$\eta$ is continuous and is trivial on the pro-$\ell$
part of $\epsilon(G_L)$, since $\ell$ is odd. In
particular it is trivial when $\omega=1$.
Reduction of characteristic polynomials gives the
permutation of the four residual characters required
by \texttt{Ok.g3'}. Its conclusion is that $\eta=1$,
contradicting construction. Thus $\rho|_N$ is
absolutely irreducible, the original \texttt{TypeI}
condition.

We prove \texttt{N3}. Let $H$ be the Zariski closure of
$\rho(N)$ in $\mathrm{GL}_4(E)$. It acts irreducibly.
The identity component of the closure of the cyclic
group generated by the inertia element found above
is a torus contained in $H^\circ$. The four weights
are distinct on it, since a ratio trivial on the identity
component would have finite image on the cyclic group.

We exclude the sum of four characters in
Lemma~\ref{lem:rh-four-regular-coset}. In that case the
four weight lines of $H^\circ$ are distinct and the
whole group $\rho(G_L)$ permutes them, because
$H^\circ$ is normal in its Zariski closure. The
permutation group $P\subset S_4$ is transitive by
irreducibility. On the kernel $Q$ of this permutation
action, the four diagonal eigenvalues in this basis
are characters. Their values are integral units in
the valuation ring of $E$: the characteristic
polynomials and their unit constant terms come from
the actual integral curve. Reduce these four
characters. Their direct sum has the same
characteristic polynomials on $Q$ as the restrictions
of the four given residual characters. The
semisimple-character theorem (or the finite-image
Brauer--Nesbitt theorem) identifies the two multisets.
Indeed the given restrictions have finite image,
and each reduced character value is one of the
finitely many prescribed values. The reduced
characters of the four weight lines are transitively
permuted by conjugation by $G_L$, whereas each of
the given residual characters is invariant under
conjugation. It follows that all four reduced
characters agree on $Q$. Thus every quotient of the
given residual characters factors through $P$.
Every cyclic quotient of a subgroup of $S_4$ has
order at most four: lift a generator of the quotient
to a permutation, whose order is at most four.
This contradicts \texttt{Ok.ord}, which supplies a
root character of order greater than four.
We have verified every assumption of
Lemma~\ref{lem:rh-four-regular-coset}.

Every open subgroup of the profinite group $N$ has
finite index. Its image's Zariski closure contains
$H^\circ$ since $H$ is a finite union of translates
of that closed subgroup. Thus the image of every
nonempty relatively open subset of $N$ is Zariski
dense in a set containing a coset of $H^\circ$.
The lemma shows that the discriminant of the
characteristic polynomial cannot vanish on that
whole image. This is exactly the required density
of separable spectra on $N$, not a density assertion
on the full Galois group in its place.

The \texttt{N1} conditions follow from the distinctness
on $T.U1(v)$: on that group $\epsilon=1$, so neither
equality of labeled characters nor equality after a
cyclotomic twist is possible. At the distinguished
place the reduction of the matrix has all four
eigenvalues equal, while the residue cyclotomic
character at Frobenius is not one. Thus conjugation
by the Frobenius matrix has eigenvalues reducing to
$1$. The cyclotomic-twisted intertwiner equation has
only the zero solution, even on the full generic
matrix algebra, proving \texttt{N4}. The polarization
matrix is the original symmetric invertible
$J_{\rm sp}R(c)$, with $J_{\rm sp}$ the fixed
symplectic form; the pairing computation is the one
in the seed-support proof. The original ordered
flag and residual inertia conditions follow from
the determinant-map construction. In particular
the character twist is a unit over $k[[z]]$ and
changes none of these conditions.

We have proved all four clauses of the original
\texttt{NiceI} predicate for the actual curve of
each ordinary determinant ring. The passage
through the level quotient changes none of the
density conditions: it is a continuous quotient
of $G_L$, with cyclotomic character descending;
the inverse image of a nonempty relatively open
subset of its cyclotomic kernel is nonempty and
relatively open in $N$.
Since $y$ kills the sum ideal, both ordinary
evaluations kill their respective ordinary Hecke
ideals. The already proved \texttt{P1core}
theorem therefore puts the respective Hecke
ideal in every prime below each ordinary curve
kernel. Pulling a prime $\mathfrak q$ of $B$
below $\ker(y|_B)$ back by each link map gives
two such primes. Hence $J\subset\mathfrak q$,
which is exactly \texttt{LocZero}.
\end{proof}

In particular the original sum ideal now satisfies
\texttt{fin}, \texttt{p1} and \texttt{seed} for
the fitted deep tower. The characteristic-zero
\texttt{pur} assertion at the remaining RH weights
and the restricted cuspidal attachment were proved
above. Global component-support propagation also
requires the local chart presentation and its
reducible-locus dimension estimate; the argument
here does not replace those by the existence of
regular spectra.
